\section{研究背景}

在传统的顶点笛卡尔坐标与单元材料标架表示下，不可伸长性通常被写为相邻离散节点之间的非线性长度约束。它们对应于离散系统中极硬的轴向模式；相比之下，杆的弯曲和扭转则对应低能量的软模式。由于这些模式被混合在同一组广义坐标中，系统既要保留细杆真实的低频弯扭响应，又不得不在求解阶段反复处理由长度约束引入的极高刚度，从而形成典型的病态系统。从表示角度看，这种病态性的根源在于原始自由度并不顺应细杆的内在运动学结构。

\subsection{相关工作与传统方法的局限}

% [Covered in Intro] 针对这些内在约束，现有方法通常采取三类思路：在原始空间中通过罚函数、拉格朗日乘子或投影策略显式施加约束；通过半内禀表示部分消除冗余变量；或者采用更彻底的紧凑表示以完全内化不可伸长性。
% [Covered in Intro]
% [Covered in Intro] 许多方法~\cite{spillmann_corde_2007,spillmann2008cosseratnet,wen_cosserat_2020}通过中心线上顶点的笛卡尔坐标来表示杆的形状，但在这种原始表示中，每个单元上的长度约束是不可避免的。此外，还需要约束四元数保持单位长度且部分轴对齐切线。为了处理这些显式内在约束，最常见的是基于惩罚的方法、拉格朗日乘子及其变体。拉格朗日乘子方法~\cite{goldenthal_efficient_2007,spillmann_inextensible_2010,deul_direct_2018}精度良好，但求解包含许多约束的 KKT 系统代价高昂。步进和投影（SAP）策略虽能强制执行解耦约束~\cite{spillmann_corde_2007}，却通常会产生物理伪影。
% [Covered in Intro]
% [Covered in Intro] 另一方面，为了消除单位四元数和对齐约束，\cite{bergou_discrete_2008,bergou_discrete_2010}等半内禀方法采用 Bishop 框架。然而其长度约束仍是非线性的，系统仍需迭代求解带有参数控制终止条件的 KKT 系统，这依然无法摆脱原空间中高刚度带来的条件数恶化与性能惩罚问题。
% [Covered in Intro]
% [Covered in Intro] 而试图彻底消除约束的全内禀紧凑表示（如 Super-Helices~\cite{bertails_super-helices_2006,bertails_linear_2009}），由于需要从杆的起点积分曲率导数以获取世界标架，会导致密集的质量矩阵与复杂的算法逻辑，甚至难以与高度稳定的隐式时间步进方案兼容。另一种基于铰接刚体链的表示策略~\cite{hadap_oriented_2006,wang_redmax_2019}虽然避免了空间积分积分困难，但当面对成百上千个微小节点时，传统的块对角预处理技术无法有效传递细长链的全局弯曲耦合。总体而言，如何在“约束消除”“代数稀疏性”和“可高效隐式积分”之间取得平衡，仍是不可伸长细杆仿真求解器的核心难题。

\subsection{本章方法与贡献}

本章的基本思想是，不再在原始笛卡尔坐标中对抗不可伸长约束，而是通过与细杆内在自由度相匹配的广义坐标重构，将这部分极硬的轴向约束直接吸收到运动学表示之中。具体而言，我们将细杆视为刚性单元链，以根顶点位置与各单元轴角变量作为状态变量，构造出一种恰好匹配内在自由度的紧凑表示。这样一来，传统表示中的长度约束、单位四元数约束与对齐约束均被自然消除，原本主导条件数的硬模式不再显式暴露于底层求解器。

在此基础上，本章进一步利用该表示下 Hessian 的特殊非零结构，提出一种新的、证明一致对称正定的混合预条件子构造方法，并设计出近线性复杂度的矩阵-向量乘法，使隐式 SQP 积分即使在大时间步长和复杂场景碰撞下，仍能实现比传统预处理框架高出一到两个数量级的性能提升与稳定求解。

总之，本工作的贡献包括：
\begin{itemize}
\item 一种紧凑表示方法，彻底消除了不可伸长 Cosserat 杆的所有内在物理硬约束，且十分契合大时间步长的隐式积分器。
\item 一种在构建和求解阶段均具有线性复杂度的有效混合预处理器，显著压低了系统条件数并加速了非线性收敛。
\item 给出了预处理器对称正定性的严格数学证明，为系统的稳定演化奠定了理论基础。
\end{itemize}
